\section{The eight-node transition and its two small modes}\label{sec:eight-node}

We apply Theorem~\ref{thm:matrix} to the compact transition constructed
in stability--degeneration Paper~2,
\cite[Appendix~A, Proposition~A.1]{wuebben-chambers}. This uses its
geometric realization and exceptional-curve data as premises. The
additional deformation argument below supplies a polarized smoothing
with the precise first-order and canonical-bundle properties required
by the analytic theorem.
This family has dual-edge length four at its nontriangular toric faces;
it is outside the length-one admissible class of CY~VI
\cite[Definition~2.1]{wuebben-cy-mirror}.

\subsection{A polarized simple smoothing of the earlier example}
Let $B=\PP^1\times\PP^1$, $L=\OO_B(1,-1)$, and
\begin{equation}\label{eq:eight-ambient}
 Z=\PP_B(\OO\oplus L)\times_B\PP_B(\OO\oplus L).
\end{equation}
Use the quotient convention for the tautological classes $\xi_1,\xi_2$,
and write $f,g$ for the base hyperplane classes. Let $\bar Z$ be the
anticanonical model, $Y\subset Z$ a general anticanonical hypersurface
chosen simultaneously nondegenerate for its toric polytope, and
$q:Y\to X_0\subset\bar Z$ the induced contraction. Set
\begin{equation}\label{eq:eight-polarization}
 \mathcal L_0=(-K_{\bar Z})|_{X_0},\qquad
 H=c_1(q^*\mathcal L_0)=(4g+2\xi_1+2\xi_2)|_Y.
\end{equation}

\begin{proposition}\label{prop:eight-family}
The threefold $X_0$ has eight ordinary double points and the displayed
map is a smooth simply connected projective small crepant resolution.
There is a proper flat projective smoothing over a disc, with trivial
relative canonical bundle and a relatively ample line bundle restricting
to the exact $\mathcal L_0$, such that every local smoothing function
has nonzero first derivative. The smooth fibres are simply connected,
and the second Betti numbers of $Y$ and $X_t$ are four and two.
\end{proposition}
\begin{proof}
The cited construction proves that $\bar Z$ has two disjoint singular
projective lines with transverse ordinary-node singularities. A general
anticanonical section meets each in four simple points. Its projective
crepant toric refinement restricts to $q$, and the resulting smooth
$Y$ has Hodge pair $(4,60)$, trivial canonical bundle and
$H^1(\OO_Y)=H^2(\OO_Y)=0$. Rationality of the nodes transfers both
vanishings to $X_0$. The model $\bar Z$ is Gorenstein Fano, so
$\mathcal L_0$ is ample and Cartier. The canonical-class calculation
in the quotient convention gives \eqref{eq:eight-polarization}, and
Cartier adjunction gives $\omega_{X_0}\simeq\OO_{X_0}$.

The four-plus-four multiplicity and the curve periods can also be read
directly from the bundle construction. Let $S_f,S_g\simeq B$ be the
sections for which both quotients are respectively onto $\OO$ and onto
$L$. Their normal bundles and anticanonical restrictions are
\[
 N_{S_f/Z}=(L^{-1})^{\oplus2},\quad N_{S_g/Z}=L^{\oplus2},\qquad
 (-K_Z)|_{S_f}=\OO(0,4),\quad (-K_Z)|_{S_g}=\OO(4,0).
\]
The restrictions of global sections are complete linear systems: the
all-$\OO$ and all-$L$ summands in
\[
 \pi_*\OO_Z(-K_Z)=\OO(0,4)\otimes\operatorname{Sym}^2(\OO\oplus L)
                              \otimes\operatorname{Sym}^2(\OO\oplus L)
\]
have independent coefficients. A general section therefore meets each
$S$ in four distinct fibres of its contracted ruling. For one such fibre
$C$, the sequence with $N_{C/S}=\OO$ and $N_{S/Z}|_C=\OO(-1)^{\oplus2}$
splits, since $H^1(\PP^1,\OO(1))=0$. The normal derivative of the
anticanonical equation is an isomorphism on the trivial summand because
the corresponding quartic zero is simple. Hence
$N_{C/Y}=\OO(-1)^{\oplus2}$. The divisor degrees in the basis
$(f,g,\xi_1,\xi_2)$ are
\begin{equation}\label{eq:eight-period-types}
 \Gamma_f=(1,0,0,0),\qquad \Gamma_g=(0,1,-1,-1).
\end{equation}
The imported global contraction calculation excludes further exceptional
loci or singularities, and its Picard-rank computation shows that these
four divisor classes span $H^2(Y;\QQ)$. Thus the displayed degrees
determine the rational exceptional homology classes.
Its full rational relation space is given by separate sum-zero
conditions on these two groups of four. In particular
$(1,1,1,-3;1,1,1,-3)$ has full support.
The first-order criterion of
\cite[Theorem~1.2 and Section~3, (3.2)]{rollenske-thomas} gives a
deformation tangent nonzero in every local node-smoothing space.
Theorem~1.6 of \cite{friedman-unobstructed} makes this tangent
unobstructed: $X_0$ is a compact Gorenstein Calabi--Yau threefold,
its isolated rational lci singularities are Du~Bois, its projective
resolution satisfies the $\partial\bar\partial$ lemma, and
$H^1(\OO_{X_0})=0$.

We verify projective effectivity and the exact polarization.
Choose $N$ so that $\mathcal L_0^N$ is very ample and
$H^1(\mathcal L_0^N)=0$. Over a small Artin extension, the
vanishing of $H^2(\OO_{X_0})$ lifts the embedding line bundle,
and $H^1(\mathcal L_0^N)=0$ lifts its coordinate sections.
The lifted morphism remains a closed immersion over the nilpotent
thickening. Thus embedded deformations map formally smoothly to
abstract deformations; equivalently the restricted Euler sequence
gives $H^1(T_{\PP^M}|_{X_0})=0$.
The corresponding Hilbert-scheme germ is smooth. Retain a lift of the
selected tangent and shrink to a neighbourhood $S$ with integral fibres.

Perform the Picard construction over this scheme $S$ before choosing
the analytic disc. After an \'etale neighbourhood, the universal family
has a section through a smooth point of $X_0$
\cite[Tag~055S]{stacks}. Its proper integral fibres give the universal
identity $\OO_T\simeq p_{T*}\OO$ under base change, where $p$ is the universal family projection.
The relative Picard scheme exists, and the section supplies a universal
Poincar\'e line bundle, normalized along that section
\cite[Theorems~2.5(2), 4.8; Exercises~3.11, 4.3 and their Appendix~A
solutions]{kleiman-picard}. Since $H^2(\OO_{X_0})=0$, the Picard
morphism is smooth at $[\mathcal L_0]$
\cite[Proposition~5.19]{kleiman-picard}. A further \'etale
neighbourhood gives a section through this point; pulling back the
Poincar\'e bundle extends precisely $\mathcal L_0$.
Ampleness is open. Both base changes have isomorphic tangent maps at
the chosen complex points. A holomorphic disc in the resulting smooth
base with the prescribed tangent therefore gives a relatively polarized
projective family and retains every nonzero local derivative. No
ramified reparametrization or replacement of $H$ by $NH$ occurs.

At a node the local equation is
$\sum z_j^2=\tau_i(t)$ with $\tau_i'(0)\ne0$, so the total
space is smooth. Properness allows the disc to be shrunk so that
the entire total space is smooth and the family is Gorenstein.
Its relative dualizing line bundle restricts to
$\omega_{X_0}\simeq\OO_{X_0}$. Since
$H^1(\omega_{X_0})=0$, cohomology and base change lifts a central
nowhere-zero generator to a relative section. Its zero set has proper
image missing the origin. Shrinking the disc removes that image and
trivializes the actual relative canonical bundle.

For simple connectivity, apply
\cite[Corollary~1.9]{batyrev-kreuzer} to the nondegenerate toric
hypersurface and the smooth projective crepant subdivision in the
construction. Its lattice quotient is zero: the relevant dual-polytope
vertices include the lattice basis $e_1,e_2,f_0,g_0$.
Thus $\pi_1(Y)=0$. The common exterior has simply connected links
$S^2\times S^3$, and the two kinds of filling retract onto $S^2$
and $S^3$. Van Kampen gives $\pi_1(X_t)=\pi_1(Y)=0$.
Finally the two independent exceptional classes give
$b_2(X_t)=4-2=2$, as also computed in the earlier proposition.
\end{proof}

These properties verify the general smoothability hypothesis of
Hein--Sun and the smoothing hypotheses of Benabida used in
Section~\ref{sec:geometry}. Locally write $\tau_i(t)=t u_i(t)$
with $u_i(0)\ne0$; a holomorphic square root of $u_i$ changes the
coordinates to the model $\sum z_j^2=t$ over the same base.
The Stenzel metric supplies the exact affine K\"ahler model. The
finite-node construction has fixed positive leading scale ratios,
and neither metric theorem imposes $b_2(X_t)=1$.
In particular the actual metrics in class $H$ satisfy the comparisons
used in Theorem~\ref{thm:matrix}.

\subsection{The quotient mass and the two coefficients}
The earlier construction calculates
\begin{equation}\label{eq:eight-intersection}
 \left(\int_YD_iD_jH\right)_{i,j}
 =\begin{pmatrix}
 0&8&4&4\\8&0&12&12\\4&12&-8&0\\4&12&0&-8
 \end{pmatrix},\qquad (D_i)=(f,g,\xi_1,\xi_2).
\end{equation}
It follows directly from
$f^2=g^2=0$, $\xi_i^2=(f-g)\xi_i$ and
$\int_Zfg\xi_1\xi_2=1$, by integrating $D_iD_j(-K_Z)^2$;
the same calculation gives $H^3=320$.
The restriction kernel has basis
\[
 k_1=g+\xi_1,\qquad k_2=g+\xi_2,
 \qquad H=2k_1+2k_2.
\]
Choose quotient lifts $f,g$. By \eqref{eq:eight-period-types}, their
period vectors have disjoint four-entry unit supports. The energy
matrix is $\mathsf E=4I_2$. The intersection blocks are
\[
 B_{KK}=\begin{pmatrix}16&24\\24&16\end{pmatrix},\quad
 B_{DK}=\begin{pmatrix}12&12\\12&12\end{pmatrix},\quad
 B_{DD}=\begin{pmatrix}0&8\\8&0\end{pmatrix}.
\]
Consequently \eqref{eq:mass-schur} gives
\begin{equation}\label{eq:eight-mass}
 \mathsf M=\frac45\begin{pmatrix}9&-1\\-1&9\end{pmatrix}.
\end{equation}
The orthogonal primitive lifts
\[
 D_-=f-g,\qquad D_+=f+g-\frac35(k_1+k_2)
\]
have squared masses $16$ and $64/5$, respectively, and both have
squared period norm $8$. Thus the generalized eigenvalues of
$(\mathsf E,\mathsf M)$ are $1/2$ and $5/8$.
Theorem~\ref{thm:matrix} proves \eqref{eq:eight-spectral-law},
including the count, eventual simplicity, type and gaps, and completes
the proof of Corollary~\ref{cor:eight-node}.

Only rational divisor data are used here. No new integral basis,
saturation or full cohomology-ring assertion is made for this family.
The application uses the earlier compact construction to derive new
information about its actual metric operators. It preserves the
distinction established there between tangent-stability walls in the
ambient fourfold, toric contraction circuits and Calabi--Yau cycle
relations.
